% GENERATED by tools/num_split.py --tex. Do not edit; rerun the tool.
{\Large\bfseries\color{acc} Marks Weightage: Numerical vs Theory}
\vspace{-6pt}{\color{acc}\rule{\textwidth}{1.1pt}}\vspace{8pt}

\begin{center}
\begin{tabularx}{\textwidth}{@{}L{0.6cm}XL{1.1cm}L{1.1cm}L{1.5cm}L{1.2cm}L{1.6cm}@{}}
\toprule
\hrow \thead{Ch} & \thead{Chapter} & \thead{Marks} & \thead{Share} & \thead{Numerical} & \thead{Theory} & \thead{Num.\,\%} \\
1 & Introduction & 159 & 8.2\% & 0 & 159 & 0\% \\
2 & Data Preprocessing & 295 & 15.2\% & 49 & 246 & 17\% \\
3 & Classification & 420 & 21.7\% & 167 & 253 & 40\% \\
4 & Association Analysis & 397 & 20.5\% & 157 & 240 & 40\% \\
5 & Cluster Analysis & 366 & 18.9\% & 135 & 231 & 37\% \\
6 & Anomaly / Fraud Detection & 158 & 8.2\% & 3 & 155 & 2\% \\
7 & Advanced Applications & 141 & 7.3\% & 8 & 133 & 6\% \\
\midrule
& \textbf{Total} & \textbf{1936} & 100\% & \textbf{519} & \textbf{1417} & \textbf{27\%} \\
\bottomrule
\end{tabularx}
\end{center}

{\small\color{sub} Marks and shares are every mark the Detailed PYQ records, recounted by \texttt{tools/num\_split.py}. A question counts as \emph{numerical} when it asks for a calculation, construction, trace or worked derivation-with-numbers of the kind the numerical companion solves; its whole marks go to that side, including a 2-mark definition asked alongside. Bold marks a chapter that is at least half numerical. Totals exceed 80 per paper because both sides of an ``or'' are counted: read these as shares, not as an audit.}
